\begin{afterword}
\summaryhead

The post asks why rationalists should care about truth and gives three answers, each of which ranks questions differently. Some people treat truth-seeking as noble, and often as a moral duty. The author is suspicious of this motive, not because the ideal is bad, but because a duty can teach bad habits of thought, as in the popular image of Spock, and people then defend their conduct instead of learning from mistakes.

Others want truth for a practical goal, such as building an airplane or finding chocolate milk. They rank questions by how much an answer could change a choice, and failure tells them which ways of thinking work. Others are curious, and rank questions by how interesting they are. Curiosity is pure, but it may not check its answers; the author says science began when people noticed that some ways of thinking let them manipulate the world. Deliberate rules of thinking are also needed. All three motives matter, but asked which is most foundational, the author answers: curiosity.

\responsehead

The post's scheme is simple and useful: three motives for seeking truth, each with its own way of ranking questions. The practical motive is the only one given a test: failure brings feedback, and a plane that falls shows that someone reasoned badly.

The case against the moral motive is thinner. Its example is Spock, a fictional character who shows a mistaken idea of rationality, all calm and more digits than the numbers deserve. He does not show that holding rationality as a duty produces the mistake, which is what the example is placed to show. The mechanism, a duty that behaves like ``an arbitrary tribal custom'' and people who ``indignantly protest'' instead of learning, comes with no case at all.

The history is asserted. What ``set humanity firmly on the path of Science'' was truth as an instrument, the post says, with no source. It adds that campfire tales satisfied curiosity and ``no one realized that anything was wrong with that.'' That is overstated if it covers the ancient Greeks the same paragraph names: Xenophanes objected to the poets' stories of the gods in the late sixth century BC.

Two turns at the end sit uneasily with what came before. First, the post calls for rules of thinking that look like norms of ``propriety,'' having blamed moral motives because people defend their errors by appeal to propriety. It does not say what keeps the new norms from that trouble. Second, it names curiosity the most foundational motive. The post's own arguments point elsewhere: curiosity may not verify its answers, campfire tales satisfied it, and the practical motive is the one credited with feedback and with science. The reason given for the choice is that the author ``just really want[s] to know.'' That is a candid personal answer. The argument of the post does not support it.

In short: a useful sorting of motives for seeking truth, whose case against the moral motive offers only a fictional example, and whose final choice of curiosity is not supported by the post's own reasons.
\end{afterword}
